\documentclass[10pt,a4paper]{article}

\usepackage[ngerman]{babel}
\usepackage{fontspec}
\setmainfont{lmroman10-regular.otf}[
  BoldFont=lmroman10-bold.otf,
  ItalicFont=lmroman10-italic.otf,
  BoldItalicFont=lmroman10-bolditalic.otf
]
\setsansfont{lmsans10-regular.otf}[
  BoldFont=lmsans10-bold.otf,
  ItalicFont=lmsans10-oblique.otf,
  BoldItalicFont=lmsans10-boldoblique.otf
]
\setmonofont{lmmono10-regular.otf}[
  Scale=0.84,
  AutoFakeBold=2,
  AutoFakeSlant=0.2
]

\usepackage[a4paper,left=19mm,right=19mm,top=18mm,bottom=19mm,headheight=14pt]{geometry}
\usepackage{microtype}
\usepackage{parskip}
\usepackage{xcolor}
\usepackage{booktabs}
\usepackage{tabularx}
\usepackage{array}
\usepackage{enumitem}
\usepackage{listings}
\usepackage{tikz}
\usetikzlibrary{arrows.meta,positioning,fit,calc,backgrounds,shapes.geometric}
\usepackage{graphicx}
\usepackage{caption}
\usepackage{float}
\usepackage{fancyhdr}
\usepackage{titlesec}
\usepackage{hyperref}
\usepackage{lastpage}

\definecolor{Primary}{HTML}{244A3B}
\definecolor{PrimaryLight}{HTML}{E9F1ED}
\definecolor{Accent}{HTML}{B46A3C}
\definecolor{AccentLight}{HTML}{F6ECE5}
\definecolor{Ink}{HTML}{1F2925}
\definecolor{Muted}{HTML}{5F6D67}
\definecolor{Line}{HTML}{C7D2CC}
\definecolor{Success}{HTML}{2F7D4B}
\definecolor{Warn}{HTML}{B7791F}
\definecolor{Failure}{HTML}{A23B3B}
\definecolor{Broker}{HTML}{E8EEF7}
\definecolor{Database}{HTML}{EEEAF7}

\newcommand{\littleGap}{\vspace{8mm}}
\newcommand{\midGap}{\vspace{12mm}}

\hypersetup{
  colorlinks=true,
  linkcolor=Primary,
  citecolor=Accent,
  urlcolor=Accent,
  pdftitle={Architektur und Umsetzung eines exemplarischen Online-Stores "Das Holzwerk"	Leading-Edge-Development},
  pdfsubject={Vertiefung Softwaretechnik FHSWF2026 - Projektbericht},
  pdfauthor={(B.Sc.) P. Förderer}
}

\pagestyle{fancy}
\fancyhf{}
\fancyhead[L]{\sffamily\small Holzwerk-OnlineStore}
\fancyhead[R]{\sffamily\small Technischer Projektbericht}
\fancyfoot[L]{\sffamily\footnotesize Stand: 31. August 2026}
\fancyfoot[R]{\sffamily\footnotesize Seite \thepage\ von \pageref{LastPage}}
\renewcommand{\headrulewidth}{0.4pt}
\renewcommand{\footrulewidth}{0pt}

\titleformat{\section}
  {\Large\sffamily\bfseries\color{Primary}}
  {\thesection}{0.7em}{}
  [\vspace{-0.25em}\color{Line}\titlerule]
\titleformat{\subsection}
  {\large\sffamily\bfseries\color{Primary}}
  {\thesubsection}{0.6em}{}
\titleformat{\subsubsection}
  {\normalsize\sffamily\bfseries\color{Ink}}
  {\thesubsubsection}{0.5em}{}
\titlespacing*{\section}{0pt}{1.0em}{0.55em}
\titlespacing*{\subsection}{0pt}{0.75em}{0.35em}
\titlespacing*{\subsubsection}{0pt}{0.55em}{0.25em}

\setlist[itemize]{leftmargin=1.25em,itemsep=1.5pt,topsep=3pt}
\setlist[enumerate]{leftmargin=1.45em,itemsep=1.5pt,topsep=3pt}
\setlength{\tabcolsep}{5pt}
\setlength{\emergencystretch}{2em}
\renewcommand{\arraystretch}{1.14}
\captionsetup{font=small,labelfont={bf,sf},textfont=it,skip=4pt}
\setcounter{tocdepth}{2}
\setcounter{secnumdepth}{3}

\newcolumntype{Y}{>{\raggedright\arraybackslash}X}
\newcolumntype{R}{>{\raggedleft\arraybackslash}p{1.45cm}}
\newcommand{\code}[1]{\texttt{\detokenize{#1}}}
\newcommand{\term}[1]{\textsf{\textbf{#1}}}
\newcommand{\good}{\textcolor{Success}{\textbf{ja}}}
\newcommand{\partialok}{\textcolor{Warn}{\textbf{teilweise}}}

\lstdefinelanguage{json}{
  basicstyle=\ttfamily\scriptsize,
  string=[s]{"}{"},
  stringstyle=\color{Accent},
  showstringspaces=false,
  breaklines=true,
  literate=
   *{0}{{{\color{Primary}0}}}{1}
    {1}{{{\color{Primary}1}}}{1}
    {2}{{{\color{Primary}2}}}{1}
    {3}{{{\color{Primary}3}}}{1}
    {4}{{{\color{Primary}4}}}{1}
    {5}{{{\color{Primary}5}}}{1}
    {6}{{{\color{Primary}6}}}{1}
    {7}{{{\color{Primary}7}}}{1}
    {8}{{{\color{Primary}8}}}{1}
    {9}{{{\color{Primary}9}}}{1}
}
\lstset{
  basicstyle=\ttfamily\scriptsize,
  frame=single,
  rulecolor=\color{Line},
  backgroundcolor=\color{PrimaryLight!45},
  keywordstyle=\color{Primary}\bfseries,
  commentstyle=\color{Muted}\itshape,
  stringstyle=\color{Accent},
  numbers=none,
  showstringspaces=false,
  breaklines=true,
  columns=fullflexible,
  aboveskip=5pt,
  belowskip=5pt,
  xleftmargin=3pt,
  xrightmargin=3pt
}

\tikzset{
  actor/.style={
    draw=Primary, fill=PrimaryLight, rounded corners=2pt,
    minimum width=19mm, minimum height=8mm, align=center,
    font=\sffamily\scriptsize\bfseries
  },
  component/.style={
    draw=Primary, fill=white, rounded corners=2pt,
    minimum width=24mm, minimum height=9mm, align=center,
    font=\sffamily\scriptsize
  },
  service/.style={
    component, fill=PrimaryLight
  },
  external/.style={
    component, draw=Accent, fill=AccentLight
  },
  infrastructure/.style={
    component, draw=Muted, fill=Broker
  },
  database/.style={
    cylinder, shape border rotate=90, aspect=0.25,
    draw=Muted, fill=Database, minimum width=23mm,
    minimum height=10mm, align=center, font=\sffamily\scriptsize
  },
  flow/.style={
    -{Latex[length=2mm]}, semithick, draw=Primary
  },
  async/.style={
    -{Latex[length=2mm]}, semithick, dashed, draw=Accent
  },
  note/.style={
    draw=Line, fill=white, rounded corners=2pt,
    align=left, font=\sffamily\tiny, inner sep=3pt
  },
  decision/.style={
    diamond, aspect=2.2, draw=Primary, fill=PrimaryLight,
    align=center, font=\sffamily\tiny, inner sep=1.2pt
  },
  state/.style={
    draw=Primary, fill=white, rounded corners=2pt,
    minimum width=24mm, minimum height=7mm, align=center,
    font=\sffamily\tiny
  },
  failedstate/.style={
    state, draw=Failure, fill=Failure!7
  },
  successstate/.style={
    state, draw=Success, fill=Success!7
  }
}

\begin{document}

\begin{titlepage}
  \thispagestyle{empty}
  \begin{tikzpicture}[remember picture,overlay]
    \fill[Primary] (current page.north west)
      rectangle ([yshift=-44mm]current page.north east);
    \fill[Accent] ([yshift=-44mm]current page.north west)
      rectangle ([yshift=-48mm]current page.north east);
    \draw[Line,line width=0.8pt]
      ([xshift=19mm,yshift=24mm]current page.south west)
      -- ([xshift=-19mm,yshift=24mm]current page.south east);
  \end{tikzpicture}

  \vspace*{8mm}
  {\color{white}\sffamily\bfseries\fontsize{15}{18}\selectfont
    Vertiefung Software Technik - SS26 - Online-Store\par}
  \vspace{22mm}

  {\sffamily\bfseries\color{Primary}\fontsize{28}{33}\selectfont
    Architektur und Umsetzung\\ 
	eines exemplarischen Online-Stores\\ 
	"Das Holzwerk" \\ 
	Leading-Edge-Development\par}
  \vspace{5mm}
  {\sffamily\color{Muted}\Large
    Projektbericht\par}

\vfill
  \begin{tabular}{ll}
	Gegenstand: 		& Entwurf, Implementierung und Qualitätssicherung eines verteilten Onlineshops\\
						& mit nachvollziehbaren Bestell- und Kompensationsabläufen\\
	Technologiestand: 	& .NET10, gRPC, RabbitMQ, PostgreSQL, OpenTelemetry, OpenAPI 3.1\\
    Dokumentstand: 		& 31. August 2026\\    
	Autor: 				&	B.Sc. P. Förderer\\
	E-Mail:				&	p.foerderer@fh-swf.de\\
\end{tabular}
\vspace{15mm}
\end{titlepage}

\pagenumbering{arabic}
\setcounter{page}{2}

\section*{Dokumentübersicht}
\addcontentsline{toc}{section}{Dokumentübersicht}

\begin{minipage}[t]{0.58\textwidth}
  \tableofcontents
\end{minipage}\hfill
\begin{minipage}[t]{0.37\textwidth}
  \vspace{11mm}
  \colorbox{PrimaryLight}{%
    \parbox{\dimexpr\linewidth-2\fboxsep\relax}{%
      \sffamily\small\textbf{\color{Primary}Kurzfassung}\par
      \vspace{3pt}
      \rmfamily\small
      Der Holzwerk-OnlineStore verbindet eine öffentliche REST-Fassade mit
      internen gRPC-Services. RabbitMQ entkoppelt Rechnungserstellung und
      Auditing; PostgreSQL speichert append-only Audit-Snapshots. Eine
      Payment-Fassade kapselt austauschbare Provider. Der ShopService
      orchestriert den Bestellprozess als SAGA und kompensiert Zahlungen und
      Reservierungen bei Folgefehlern.
    }}
\end{minipage}

\vspace{7mm}
\begin{tabularx}{\textwidth}{@{}p{31mm}Y@{}}
  \toprule
  \textbf{\sffamily Merkmal} & \textbf{\sffamily Ausprägung}\\
  \midrule
  System- & Browserzugriffe ausschließlich über den StoreProxy; interne
    Dienste über gRPC und RabbitMQ\\
  Qualitätsfokus & Erweiterbarkeit, Nachvollziehbarkeit, Fehlertransparenz,
    reproduzierbarer lokaler Betrieb\\
  Persistenz & PostgreSQL für Audit-Snapshots; JSON-Dateien für Lager- und
    Rechnungsdaten der Demonstrationsumgebung\\
  Nachweis & 28 automatisierte Tests sowie End-to-End-Prüfungen der vier
    deterministischen Präsentationsszenarien\\
  \bottomrule
\end{tabularx}

\newpage
\section{Einleitung}
\littleGap
\subsection{Problemstellung}
Ein Bestellprozess wirkt aus Sicht des Kunden wie eine einzelne Aktion. Intern
müssen jedoch Bestand, Zahlung, Rechnung und Protokollierung koordiniert werden.
In einem verteilten System können diese Teilschritte unabhängig voneinander
fehlschlagen oder zu spät antworten. Eine lokal erfolgreiche Zahlung bedeutet
beispielsweise noch nicht, dass der Lagerbestand endgültig ausgebucht oder eine
Rechnung erstellt wurde. Klassische Datenbanktransaktionen über alle
Systemgrenzen hinweg wären technisch aufwendig, würden die Komponenten eng
koppeln und passen nicht zu einem ereignisbasierten Rechnungs- und Auditingpfad.

Das Projekt adressiert daher drei miteinander verbundene Fragen:

\begin{itemize}
  \item Wie lässt sich der Bestellprozess in fachlich getrennte Services
    zerlegen, ohne die Konsistenzanforderungen zu verletzen?
  \item Wie werden Fehler und Timeouts so behandelt, dass Kunde und Betreiber
    einen eindeutigen Zustand erkennen können?
  \item Wie bleibt die Historie eines Vorgangs über synchrone Aufrufe,
    Nachrichten und Kompensationen hinweg nachvollziehbar?
\end{itemize}

\midGap

\subsection{Ziele}

Das Ergebnis soll eine lokal ausführbare Referenzimplementierung eines Onlineshops sein. Im Vordergrund stehen Architektur und Verhalten. Da es sich hier um ein exemplarisches Lernbeispiel handelt, wurden einige Freiheiten, wie zB. Payment-Provider Implementierung als Stub's oder ein generiertes Fronted, bereits in der Voraussetzung gewährt.\\
Daraus ergeben sich folgende Ziele:

\begin{enumerate}
  \item Eine öffentliche, ressourcenorientierte REST-API mit OpenAPI-3.1-
    Beschreibung und einheitlichen Problem-Details nach RFC~9457
  \item Getrennte Verantwortlichkeiten für Orchestrierung, Lager, Zahlung,
    Rechnung und Auditing
  \item Eine konfigurierbare Payment-Fassade mit mehreren Providern und einem
    Erweiterungspunkt nach dem Open/Closed Principle
  \item Ein expliziter SAGA-Ablauf mit definierten terminalen Zuständen und
    nachvollziehbaren Kompensationen
  \item Eine unveränderliche, korrelationsbezogene Audit-Historie in
    PostgreSQL sowie ergänzende technische Logs und Telemetrie
  \item Reproduzierbarer Betrieb über PowerShell-Skripte und eine
    vorführbare Fault Injection pro Bestellung
\end{enumerate}

\midGap

\subsection{Abgrenzung}

Die Implementierung ist eine Entwicklungs- und Demonstrationsumgebung. PayPal,
Stripe und DemoPay sind lokale Stubs und bewegen kein echtes Geld. Die
Rechnungsmail wird standardmäßig in eine Pickup-Outbox geschrieben; ein
produktiver SMTP-Versand, Webhooks externer Provider, Authentifizierung,
Autorisierung, Hochverfügbarkeit und Backupkonzepte sind
nicht oder nur zum Teil, Gegenstand des Projekts. Nur Audit-Snapshots benötigen relationale
Persistenz nach Voraussetzung. Lager- und Rechnungsdaten bleiben bewusst dateibasiert. Diese
Abgrenzung verhindert, dass Infrastrukturthemen die Untersuchung der
Architektur- \& Entwurfsmuster überlagern.\\
Auch das Frontend (index.html, css- \& js-Files) durfte vernachlässigt werden und wurde durch 
ChatGPT's-Codex-Plattform\cite{cgpt} generiert, um den Fokus nicht auch auf Webdesign und JavaScript zu legen.

\newpage

\section{Grundlagen}

\littleGap

\subsection{Microservices und Gateway}


Die Services orientieren sich an fachlichen Verantwortlichkeiten. Der
StoreProxy ist der einzige öffentliche Einstiegspunkt und leitet explizit
freigegebene HTTP-Ressourcen an den ShopService weiter. Dadurch kennt der
Browser weder interne Ports noch gRPC-Verträge. Der ShopService besitzt die
Prozesssicht, während Warehouse-, Billing-, Invoice- und AuditService ihre
jeweiligen Fachfunktionen kapseln. Diese Trennung unterstützt unabhängige
Fehlerbehandlung und vermeidet direkte Zugriffe der Website auf interne
Komponenten. Während die YARP-Middleware hier, zentrales Routing und LoadBalancing
übernimmt.

Die öffentliche API folgt dem Ressourcenstil: Bestellungen werden mit
\texttt{POST} \path|/api/orders| erzeugt, Rechnungen über
\texttt{GET} \path|/api/invoices/{invoiceId}/pdf| abgerufen und
Audit-Historien über \texttt{GET}
\path|/api/order-audits/{correlationId}/snapshots| gelesen. Aktionsnamen
wie \code{/chargeOrder} oder \code{/getAudit} werden vermieden. OpenAPI 3.1
beschreibt Requests, Antworten, Statuscodes und Fehlerfälle. Problem-Details
nach RFC~9457 liefern neben \code{type}, \code{title}, \code{status},
\code{detail} und \code{instance} stets eine \code{correlationID}
\cite{rfc9457,openapi}.

\midGap

\subsection{Fassade und Provider-Strategien}

Das Fassadenmuster bietet einem Client eine vereinfachte Schnittstelle zu einem
komplexeren Teilsystem \cite{gamma}. Im Projekt ist die
\code{PaymentFacade} der einzige Zugriffspunkt auf konkrete
\code{IPaymentProvider}-Implementierungen. Sie vereinheitlicht die Operationen
\code{ChargeAsync}, \code{RefundAsync} und \code{GetStatusAsync}, validiert die
Providerwahl und behandelt Zeitüberschreitungen an einer Stelle.

Die konkreten Adapter bilden quasi austauschbare Strategien im Sinne des Entwurfsmusters, 
PayPal und Stripe sind initial aktiviert; DemoPay bleibt registriert, ist jedoch
deaktiviert. Die Website trifft keine Vorauswahl. Der Kunde wählt den Provider
für genau eine Bestellung, und der BillingService prüft diese Wahl erneut gegen
die Konfiguration. Diese Konfiguration kann auch über die Datei: \emph{appsettings.json} des 
Projekts \emph{BillingService}, zB. auch als Umgebungsvariable gesetzt werden. Zur Kenntlichmachung 
der zusätzlichen Konfigurationsebene wurde DemoPay beispielhaft deaktiviert. Neue Adapter werden per 
Tyediscovery registriert. Deshalb müssen weder Fassade noch Startcode verändert werden, solange der neue Typ die
Schnittstelle \code{IPaymentProvider} implementiert.

\midGap

\subsection{SAGA statt verteilter ACID-Transaktion}

Das SAGA-Pattern \emph{(dt. Sage, Erzählung)} zerlegt einen Geschäftsprozess in lokale Transaktionen und ordnet
erfolgreichen Schritten passende Gegenaktionen zu \cite{richardson}. Im
vorliegenden System orchestriert der ShopService diese Reihenfolge. Eine
Bestandsreservierung wird bei abgelehnter Zahlung freigegeben. Ist die Zahlung
bereits erfolgreich und scheitern anschließend Event-Publikation oder
Lager-Commit, wird zuerst ein Refund über die Fassade versucht und danach die
Reservierung storniert.\\
Kompensation ist nicht mit einem technischen Rollback gleichzusetzen. Sie ist
eine neue fachliche Operation, kann selbst fehlschlagen und muss deshalb einen
eigenen Audit-Zustand besitzen. Der Endzustand
\code{ROLLBACK_COMPLETED} bedeutet, dass alle vorgesehenen Gegenaktionen
erfolgreich waren, wohingegen \code{ROLLBACK_FAILED} betriebliche Probleme signalisiert.\\
\\
In Abgrenzung zu typischen ACID-Transaktionen \emph{(Akronym aus: Atomicity, Consistency, Isolation, Durability)} wie sie auch gerne bei Datenbanktransaktionen verwendet werden, kann innerhalb verteilter nebenläufiger Systeme keine \emph{strenge} Konsistenz erzwungen werden, daher bietet das SAGA-Pattern eine abgeschwächte Form, der "`eventuellen Konsistenz"'.

\midGap

\subsection{Event Sourcing Light und Beobachtbarkeit}

Vollständiges Event Sourcing rekonstruiert den aktuellen Aggregatzustand aus
dem Ereignisstrom und verwendet Ereignisse als primäre Persistenz
\cite{fowler}. Der OnlineStore verfolgt bewusst eine leichtere Variante:
fachlich relevante Zustandsübergänge werden als append-only Audit-Snapshots
gespeichert, die operativen Modelle von Lager, Zahlung und Rechnung werden
jedoch nicht daraus wiederhergestellt. Der Begriff \emph{Event Sourcing Light}
bezeichnet hier somit eine unveränderliche, geordnete Historie, nicht ein
vollständiges Event-Sourcing-System.\\
Auditing und technisches Monitoring bleiben getrennt. Audit-Snapshots erklären,
\emph{was fachlich geschah}; JSONL-Logs, Traces und Metriken erklären,
\emph{wie sich die Software technisch verhielt}. Eine durchgängige
Correlation-ID verbindet HTTP-Request, gRPC-Aufrufe, RabbitMQ-Nachrichten,
Logeinträge und Datenbankzeilen.
\vspace{15mm}

\subsection{Synchron und asynchron}
Bestandsreservierung, Zahlung und Lager-Commit beeinflussen die unmittelbare
Bestellantwort und werden synchron per gRPC ausgeführt. Rechnungserstellung und
Auditing sind zeitlich entkoppelbar und laufen über RabbitMQ. Durable Queues,
persistente Nachrichten, Publisher-Confirms, manuelle Acknowledgements und
Dead-Letter-Queues ermöglichen eine At-least-once-Verarbeitung
\cite{rabbitmq}. Daraus folgt Eventual Consistency: Eine Bestellung kann bereits
abgeschlossen sein, während Rechnung oder Audit-Snapshot noch verarbeitet
werden.

\midGap

\section{Systemarchitektur}
\littleGap

\subsection{Komponenten und Verantwortlichkeiten}

\begin{table}[H]
\centering
\small
\begin{tabularx}{\textwidth}{@{}p{30mm}p{14mm}Y@{}}
  \toprule
  \textbf{\sffamily Komponente} & \textbf{\sffamily Port} &
    \textbf{\sffamily Verantwortung}\\
  \midrule
  StoreProxy & 6680 & Öffentlicher Einstieg, Website, YARP-Routing,
    Rate Limits und Proxy-Fehlerbehandlung\\
  StoreBackend & 6681 & Dateibasierte Produkt-, Bestands- und
    Reservierungspersistenz\\
  ShopService & 6682 & REST-Fassade, Readiness-Prüfung und
    SAGA-Orchestrierung\\
  WarehouseService & 6683 & Katalog, Reservieren, Ausbuchen und Freigeben\\
  BillingService & 6684 & Payment-Fassade, Providerwahl, Charge, Refund und
    Zahlungsereignis\\
  InvoiceService & 6685 & Asynchrone Rechnung, PDF und E-Mail-Pickup-Outbox\\
  AuditService & 6686 & Konsum, Validierung, Persistenz und Abfrage von
    Audit-Snapshots\\
  OTel Collector & 6687 & Zentrale Annahme technischer Logs, Traces und
    Metriken\\
  PostgreSQL & 6688 & Append-only Persistenz der Audit-Historie\\
  RabbitMQ & 5672 & Transport von Zahlungs- und Audit-Ereignissen\\
  \bottomrule
\end{tabularx}
\caption{Laufzeitkomponenten der lokalen Umgebung}
\end{table}

\begin{figure}[H]
\centering
\resizebox{0.98\textwidth}{!}{%
\begin{tikzpicture}
  \node[actor] (customer) at (0,0) {Kunde\\Browser};
  \node[service] (proxy) at (2.8,0) {StoreProxy\\REST + Website};
  \node[service] (shop) at (5.6,0) {ShopService\\Orchestrator};
  \node[service] (billing) at (8.5,0) {BillingService\\Payment-Fassade};
  \node[external] (providers) at (11.4,0) {Provider-Stubs\\PayPal / Stripe};

  \node[service] (warehouse) at (8.5,1.55)
    {WarehouseService\\Bestand};
  \node[component] (backend) at (11.4,1.55)
    {StoreBackend\\JSON};
  \node[infrastructure] (rabbit) at (8.5,-1.55)
    {RabbitMQ\\Topic-Exchanges};
  \node[service] (invoice) at (11.4,-1.55)
    {InvoiceService\\PDF + Outbox};
  \node[service] (audit) at (5.6,-3.05)
    {AuditService\\Snapshot-Owner};
  \node[database] (postgres) at (8.5,-3.05)
    {PostgreSQL\\audit\_snapshots};
  \node[infrastructure] (otel) at (11.4,-3.05)
    {OpenTelemetry\\Logs + Traces};

  \draw[flow] (customer) -- node[above,font=\sffamily\tiny]{HTTP} (proxy);
  \draw[flow] (proxy) -- node[above,font=\sffamily\tiny]{YARP} (shop);
  \draw[flow] (shop.north) |- node[pos=0.72,above,font=\sffamily\tiny]{gRPC}
    (warehouse.west);
  \draw[flow] (warehouse) -- node[above,font=\sffamily\tiny]{gRPC} (backend);
  \draw[flow] (shop) -- node[above,font=\sffamily\tiny]{gRPC} (billing);
  \draw[flow] (billing) -- node[above,font=\sffamily\tiny]{in-process} (providers);
  \draw[async] (billing) -- node[right,font=\sffamily\tiny]{Events} (rabbit);
  \draw[async] (rabbit) -- node[above,font=\sffamily\tiny]{payment.succeeded} (invoice);
  \draw[async] (rabbit.west) -| node[pos=0.72,left,font=\sffamily\tiny]{audit.\#}
    (audit.north);
  \draw[flow] (audit) -- node[above,font=\sffamily\tiny]{EF Core / Npgsql} (postgres);

  \node[note,below=3mm of customer,anchor=north west] (boundary)
    {Öffentlich: nur StoreProxy\\Intern: gRPC + RabbitMQ};
  \node[note,below=3mm of otel] {
    Alle .NET-Services exportieren\\
    technische Telemetrie per OTLP.};
\end{tikzpicture}}
\caption{Systemkontext und zentrale Kommunikationsbeziehungen}
\label{fig:architecture}
\end{figure}
\newpage  	



	 	 	
\subsection{Interaktionen im Happy Path}

Der erfolgreiche Ablauf folgt sieben fachlichen Schritten. Die Correlation-ID
wird beim Eingang erzeugt und zugleich als \code{orderId} verwendet:

\begin{enumerate}
  \item Der StoreProxy leitet \code{POST /api/orders} weiter; der ShopService
    prüft Betriebsbereitschaft und Request
  \item Der WarehouseService reserviert die Artikel persistent im
    StoreBackend
  \item Der BillingService delegiert die Zahlung über die Payment-Fassade an
    den ausgewählten Provider
  \item Der Provider bestätigt; der BillingService protokolliert das Ergebnis
  \item Der BillingService publiziert \code{payment.succeeded}; der
    InvoiceService erstellt die Rechnung asynchron
  \item Der WarehouseService bucht die reservierten Waren endgültig aus
  \item Der ShopService schreibt \code{ORDER_COMPLETED} und antwortet mit
    \code{201 Created}. Er wartet nicht auf das fertige PDF
\end{enumerate}

\littleGap

\begin{figure}[H]
\centering
\resizebox{\textwidth}{!}{%
\begin{tikzpicture}[x=2.8cm,y=-0.72cm,font=\sffamily\tiny]
  \foreach \x/\name in {
    0/Kunde + Proxy,
    1/ShopService,
    2/Warehouse,
    3/Billing + Fassade,
    4/RabbitMQ,
    5/InvoiceService} {
      \node[component,minimum width=23mm] (p\x) at (\x,0) {\name};
      \draw[Line,dashed] (\x,0.35) -- (\x,8.0);
  }

  \draw[flow] (0,1.0) -- node[above]{1: POST /api/orders} (1,1.0);
  \draw[flow] (1,2.0) -- node[above]{2: ReserveCart} (2,2.0);
  \draw[flow] (2,2.55) -- node[above]{reserviert} (1,2.55);
  \draw[flow] (1,3.35) -- node[above]{3: ProcessPayment} (3,3.35);
  \draw[flow] (3,3.9) -- node[above]{4: Transaktion erfolgreich} (1,3.9);
  \draw[async] (3,4.75) -- node[above]{5: payment.succeeded} (4,4.75);
  \draw[async] (4,5.35) -- node[above]{asynchrone Verarbeitung} (5,5.35);
  \draw[flow] (1,6.15) -- node[above]{6: CommitCart} (2,6.15);
  \draw[flow] (2,6.7) -- node[above]{Bestand ausgebucht} (1,6.7);
  \draw[flow] (1,7.5) -- node[above]{7: 201 / COMPLETED} (0,7.5);

  \node[note,anchor=west] at (4.1,6.2)
    {Rechnung und E-Mail-Outbox\\laufen unabhängig weiter};
\end{tikzpicture}}
\caption{Happy-Path-Sequenz der implementierten Bestelllogik}
\label{fig:happypath}
\end{figure}

\midGap

\subsection{Schnittstellen und Betriebsmodell}

Die internen Verträge der gRPC-Services, liegen als typische Protobuf-Dateien vor. Der öffentliche
StoreProxy kennt diese Verträge aber nicht, sondern routet ausschließlich HTTP zum
ShopService. YARP prüft dessen Health-Endpunkt aktiv, setzt eine kanonische
\code{X-Correlation-ID} und übersetzt Weiterleitungsfehler in
Problem-Details. Rate Limits begrenzen insbesondere Bestellungen auf zehn
Requests pro Client und Minute. Der Checkout selbst wird wegen möglicher
Zahlungswirkungen nicht automatisch wiederholt.\\
Die Verwendung der YARP-Middleware \cite{yarp} ermöglicht völlig nativ, die
Webanwendung zu routen (per ReverseProxy und Layer7-Routing), LoadBalancing zu nutzen,
Health-Checks durchzuführen sowie die SSL/TLS-Zertifikatsverwaltung zu übernehmen (HTTP2.0).

Das Startskript \emph{Start-VSTOnlineStore.ps1} verwaltet PostgreSQL, OpenTelemetry Collector und alle
.NET-Services komfortabel und in definierter Reihenfolge. PostgreSQL wird beim Stop zuletzt
kontrolliert beendet. RabbitMQ bleibt eine externe aber notwendige Komponente und muss installiert oder per
Docker gehostet werden.\\
Einzelne Komponenten können unabhängig gestartet und gestoppt werden; dadurch lassen
sich Ausfälle und Wiederanlauf gezielt untersuchen.\\
Website, Services und Komponenten können aber auch einzeln populiert und zur Ausführung gebracht werden,
davon wurde aber unter dem Gesichtspunkt der Präsentationsfähigkeit (Live-Demo) Abstand genommen, aber 
dennoch über den GIT-Ordner \emph{ReleaseBuild}, inklusive Anleitung für das manuelle Setup, berücksichtig.

\newpage

\section{Implementierung}
\littleGap

\subsection{Payment-Fassade}

Der gRPC-Service \code{BillingOperationsGrpcService} ist die fachliche
Servicegrenze, der konkrete Provider wird ausschließlich hinter der Fassade
aufgerufen. Die Schnittstelle bietet für alle Adapter dieselben drei
Operationen:

\begin{lstlisting}[language={[Sharp]C},caption={Vereinfachter Providervertrag}]
public interface IPaymentProvider
{
    Task<PaymentChargeResult> ChargeAsync(
        Guid orderId, long amountInCents, string currency,
        CancellationToken cancellationToken);

    Task<PaymentRefundResult> RefundAsync(automa
        string transactionId, long amountInCents,
        CancellationToken cancellationToken);

    Task<PaymentStatusResult> GetStatusAsync(
        string transactionId, CancellationToken cancellationToken);
}
\end{lstlisting}

\midGap

\begin{figure}[H]
\centering
\resizebox{0.96\textwidth}{!}{%
\begin{tikzpicture}[node distance=7mm and 10mm]
  \node[service] (grpc) {BillingOperations\\gRPC-Service};
  \node[service,right=of grpc] (facade) {PaymentFacade\\Validierung + Timeout};
  \node[component,right=of facade] (contract) {IPaymentProvider\\Charge / Refund / Status};

  \node[external,above right=7mm and 10mm of contract] (paypal)
    {PayPalPaymentProvider\\aktiv};
  \node[external,right=of contract] (stripe)
    {StripePaymentProvider\\aktiv};
  \node[external,below right=7mm and 10mm of contract,
    draw=Muted,fill=black!4,text=Muted] (demo)
    {SimulatedProvider\\deaktiviert};
  \node[external,below=17mm of facade,dashed] (new)
    {Neuer Adapter\\z.B. Klarna};

  \draw[flow] (grpc) -- node[above,font=\sffamily\tiny]{nur über Fassade} (facade);
  \draw[flow] (facade) -- node[above,font=\sffamily\tiny]{einheitlicher Vertrag} (contract);
  \draw[flow] (contract) |- (paypal);
  \draw[flow] (contract) -- (stripe);
  \draw[flow] (contract) |- (demo);
  \draw[async] (new) -| node[pos=0.25,below,font=\sffamily\tiny]
    {automatische Typentdeckung} (contract);
  \node[note,below=4mm of new] {
    OCP: neuer Typ implementiert IPaymentProvider;\\
    Fassade, gRPC-Vertrag und ShopService bleiben unverändert.};
\end{tikzpicture}}
\caption{Payment-Fassade und Erweiterung um einen Provider}
\label{fig:facade}
\end{figure}

\littleGap

\code{EnabledProviderKeys} steuert, welche registrierten Adapter für neue
Bestellungen zulässig sind. \code{ActiveProviderKey} dient nur als
serverseitiger Fallback für ältere oder interne Aufrufer. Ein Timeout wird
zentral mit einem verknüpften CancellationToken umgesetzt und für jeden
Provider in dieselbe Fehlersemantik übersetzt.

Für Refund und Status merkt sich die Fassade den Provider einer Transaktion in
einem nebenläufigen Dictionary. Dadurch bleibt der Aufrufer providerneutral.
Diese Zuordnung ist gegenwärtig prozesslokal; ihre Persistierung ist eine
bewusste offene Anforderung und wird in Kapitel~\ref{sec:reflection}
kritisch bewertet.

\midGap

\newpage

\subsection{AuditSnapshot-Schema}
Jeder fachliche Zustand wird als \code{AuditSnapshot} abgebildet. Die
Anwendung erzeugt zunächst einen Entwurf. Der AuditService ergänzt
\code{previousEventID} und eine globale \code{sequenceNumber} atomar:

\littleGap

\begin{lstlisting}[language=json,caption={Logisches AuditSnapshot-Schema}]
{
  "eventID": "uuid",
  "correlationID": "uuid",
  "eventType": "PAYMENT",
  "responsibleService": "ShopService",
  "timestamp": "2026-08-25T10:15:22.184Z",
  "payload": {
    "phase": "PAYMENT_REFUND_REQUESTED",
    "orderStatus": "IN_PROGRESS"
  },
  "previousEventID": "uuid | null",
  "actor": "ShopService",
  "statusCode": "COMPENSATING",
  "sequenceNumber": 4711
}
\end{lstlisting}

\littleGap

PostgreSQL verwendet native \code{uuid}- und
\code{timestamp with time zone}-Typen; der Payload liegt als \code{jsonb} vor.
Primär-, Fremd-, Unique- und Check-Constraints sichern Struktur und Verkettung.
Ein \emph{transaktionsgebundener Advisory Lock} serialisiert nur Ereignisse derselben
Correlation-ID, und vermeidet gleichzeitig die Sperrung von Tabellen. Datenbanktrigger 
lehnen \code{UPDATE}, \code{DELETE} und
\code{TRUNCATE} ab. Damit wird Append-only nicht nur als
Anwendungskonvention, sondern auch in der Datenbank durchgesetzt
\cite{postgres}.

\littleGap

RabbitMQ kann Nachrichten im At-least-once-Modell erneut zustellen. Die
\code{eventID} wirkt deshalb als Idempotenzschlüssel. Ein bereits vorhandenes,
inhaltlich identisches Ereignis gilt als erfolgreiche Wiederholung; ein
widersprüchliches Duplikat wird abgelehnt. Nur der AuditService besitzt das
Schema. Andere Services publizieren Ereignisse und lesen Historien über die
interne gRPC-Schnittstelle des AuditService.

\littleGap

\subsection{Kompensationslogik}

\begin{figure}[H]
\centering
\resizebox{0.98\textwidth}{!}{%
\begin{tikzpicture}
  \node[state] (start) at (0,0) {Bestellung validiert};
  \node[decision] (reserve) at (2.25,0) {Reserve\\erfolgreich?};
  \node[decision] (charge) at (4.45,0) {Charge\\erfolgreich?};
  \node[decision] (publish) at (6.65,0) {Event\\publiziert?};
  \node[decision] (commit) at (8.85,0) {Commit\\erfolgreich?};
  \node[successstate] (complete) at (11.05,0) {COMPLETED};

  \node[failedstate] (stockfail) at (2.05,-1.55) {OUT\_OF\_STOCK};
  \node[state] (release1) at (4.55,-1.55)
    {STOCK\_RELEASE\\REQUESTED};
  \node[failedstate] (payfail) at (4.55,-2.85) {PAYMENT\_FAILED};

  \node[state] (refund) at (7.75,-1.55)
    {PAYMENT\_REFUND\\REQUESTED};
  \node[state] (release2) at (7.75,-2.85)
    {STOCK\_RELEASE\\REQUESTED};
  \node[failedstate] (rollback) at (7.75,-4.15)
    {ROLLBACK\_COMPLETED / FAILED};

  \node[note] (invoice) at (10.7,1.5)
    {Nach durable Event:\\
     3 Invoice-Versuche, dann DLQ + Auditfehler;\\
     keine automatische Zahlungserstattung};

  \draw[flow] (start) -- (reserve);
  \draw[flow] (reserve) -- node[above,font=\sffamily\tiny]{ja} (charge);
  \draw[flow] (charge) -- node[above,font=\sffamily\tiny]{ja} (publish);
  \draw[flow] (publish) -- node[above,font=\sffamily\tiny]{ja} (commit);
  \draw[flow] (commit) -- node[above,font=\sffamily\tiny]{ja} (complete);

  \draw[flow] (reserve) -- node[right,font=\sffamily\tiny]{nein} (stockfail);
  \draw[flow] (charge) -- node[right,font=\sffamily\tiny]{nein} (release1);
  \draw[flow] (release1) -- (payfail);
  \draw[flow] (publish) -- node[right,font=\sffamily\tiny]{nein} (refund);
  \draw[flow] (commit.south) |- node[pos=0.25,right,font=\sffamily\tiny]{nein}
    (refund.north);
  \draw[flow] (refund) -- (release2);
  \draw[flow] (release2) -- (rollback);
  \draw[async] (publish.north) |- (invoice.west);
\end{tikzpicture}}
\caption{SAGA-Entscheidungen und fachliche Gegenaktionen}
\label{fig:saga}
\end{figure}

\midGap

Jeder Gegenlauf beginnt mit einem eigenen Snapshot im Status
\code{COMPENSATING}, der Erfolgsfall wird mit \code{COMPENSATED} und ein Fehler mit
\code{FAILURE} gespeichert. Eine Zahlungsablehnung führt nur zur
Lagerfreigabe, weil noch keine erfolgreiche Zahlung existiert. Nach erfolgreicher
Zahlung werden bei fehlgeschlagener Event-Publikation oder fehlgeschlagenem
Commit, Refund und Lagerfreigabe, nacheinander ausgeführt. Ein technisch nicht
bestätigter Commit wird genau einmal \emph{idempotent} wiederholt, bevor die
Kompensation beginnt.\\
Ist \code{payment.succeeded} hingegen bereits dauerhaft im Broker angenommen,
gehört die Rechnungserstellung denoch zur asynchronen Verarbeitung und wird bei Fehlschlag, maximal dreimal
nacheinander ausgeführt. Die Phasen \path|INVOICE_RETRY_SCHEDULED|,
\path|INVOICE_RETRY_EXHAUSTED| und
\path|INVOICE_PROCESSING_FAILED| werden auditiert; anschließend gelangt die
Nachricht in die DLQ \emph{(engl. Dead Letter Queue)}. Die Zahlung wird in diesem Fall nicht rückgängig gemacht,
weil der Bestellabschluss bereits fachlich bestätigt und die Nachricht
dauerhaft vorhanden ist.

\midGap

\subsection{Fault Injection und Präsentationsmodus}

Der Schalter \code{-PresentationMode} aktiviert vier deterministische,
bestellungsbezogene Szenarien. Die Website zeigt sie im Warenkorb an und
überträgt den Schlüssel über REST, gRPC und RabbitMQ. Im normalen
Start ist der Modus deaktiviert; eine unbeabsichtigte globale Störung wird damit
vermieden.

\begin{table}[H]
\centering
\small
\begin{tabularx}{\textwidth}{@{}p{34mm}Yp{34mm}@{}}
  \toprule
  \textbf{\sffamily Szenario} & \textbf{\sffamily Erwarteter Ablauf} &
    \textbf{\sffamily Endzustand}\\
  \midrule
  Zahlung abgelehnt & Reservierung freigeben; kein Rechnungsevent &
    \code{PAYMENT_FAILED}\\
  Bestand fehlt & Nach Reserve-Ablehnung keine Zahlung und keine Rechnung &
    \code{OUT_OF_STOCK}\\
  Invoice nicht verfügbar & Drei Versuche, Fehler-Snapshots und DLQ; kein
    Refund & Bestellung \code{COMPLETED}, Invoice fehlerhaft\\
  Warehouse-Commit fehlerhaft & Refund über Fassade, danach Reservierung wieder freigeben & \code{ROLLBACK_COMPLETED}\\
  \bottomrule
\end{tabularx}
\caption{Deterministische Vorführszenarien}
\end{table}

Mit der Aktion \code{PresentationSnapshots} werden alle Snapshots einer
Correlation-ID per ShopService über die öffentliche Audit-Ressource gelesen 
\emph{(GET http://localhost:6680/api/order-audits/{correlationId}/snapshots)}, 
formatiert in \code{Logs/Presentation} abgelegt und in Notepad++ geöffnet. Damit zeigt die
Präsentation nicht nur die sichtbare Fehlerreaktion, sondern auch den
internen Kompensationslauf.

\midGap

\section{Qualitätssicherung}

\littleGap

\subsection{Teststrategie}

Die Testpyramide konzentriert schnelle Unit-Tests auf fachliche Regeln und
Integrationstests auf die Orchestrierung:

\begin{itemize}
  \item \term{Payment-Fassade:} Erfolg, Ablehnung, Timeout je Provider,
    Konfiguration, Providerwahl, deaktivierter Adapter, Refund, Status und
    automatische Adaptererkennung.
  \item \term{StoreBackend:} persistente und idempotente Reservierung,
    Commit sowie einmalige Freigabe.
  \item \term{Shop-Integration:} Happy Path, Bestands- und Zahlungsfehler,
    Readiness, Downstream-Ausfall, Timeouts, Event-Publikationsfehler,
    Commit-Retry, Kompensation, Providerpropagation und Präsentationsmodus.
  \item \term{End-to-End:} Start des realen lokalen Stacks, Checkout über die
    Website und Abfrage der in PostgreSQL gespeicherten Snapshot-Kette.
\end{itemize}

Die Billing-Unit-Tests verwenden kontrollierte Provider-Doubles. Die
Shop-Integrationstests ersetzen gRPC-Clients und Audit-Recorder gezielt, sodass
Aufrufreihenfolge, Nichtaufrufe und terminale Zustände beobachtbar sind.
Technische Fehler werden nicht nur anhand des HTTP-Status bewertet: Die Tests
prüfen zusätzlich Refund-, Release- und Audit-Aufrufe. Dadurch kann eine
oberflächlich korrekte Fehlerantwort keinen fehlenden SAGA-Schritt verdecken.

\midGap

\subsection{Automatisierte Testergebnisse}

Der abschließende Release-Build wurde mit .NET10-64bit erfolgreich ausgeführt.

\begin{table}[H]
\centering
\begin{tabularx}{0.84\textwidth}{@{}Y>{\centering\arraybackslash}p{14mm}
  >{\centering\arraybackslash}p{14mm}>{\centering\arraybackslash}p{14mm}@{}}
  \toprule
  \textbf{\sffamily Testprojekt} &
    \textbf{\sffamily Erfolg} &
    \textbf{\sffamily Fehler} &
    \textbf{\sffamily Gesamt}\\
  \midrule
  StoreBackend.UnitTests & 2 & 0 & 2\\
  BillingService.UnitTests & 10 & 0 & 10\\
  ShopService.IntegrationTests & 16 & 0 & 16\\
  \midrule
  \textbf{Gesamt} & \textbf{28} & \textbf{0} & \textbf{28}\\
  \bottomrule
\end{tabularx}
\caption{Testergebnis des Release-Builds}
\end{table}

\littleGap

Nach Abschluss wurde die Produktversion für alle Komponenten auf Version 1.0 gesetzt. Ein fehlerfreier Release-Build populiert außerdem alle notwendigen Komponenten (außer RabbitMQ) im Verzeichnis \code{..\MASTER\VST_OnlineStore\ReleaseBuild}, ein Zip-File mit dessen Inhalt wurde unter GIT ebenfalls deployed (\code{VST_OnlineStore-1.0.0-win-x64.zip}).

\subsection{Nachweise der Fehlerszenarien}

\begin{table}[H]
\centering
\small
\begin{tabularx}{\textwidth}{@{}p{34mm}p{34mm}Y@{}}
  \toprule
  \textbf{\sffamily Fehlerfall} & \textbf{\sffamily Sichtbares Ergebnis} &
    \textbf{\sffamily Nachgewiesene interne Reaktion}\\
  \midrule
  Zahlung abgelehnt & HTTP 422,\newline
    \code{PAYMENT_FAILED} & Separate Release-Anforderung,
    persistente Freigabe und terminaler Snapshot; kein Commit\\
  Unzureichender Bestand & HTTP 409,\newline
    \code{OUT_OF_STOCK} & Abbruch nach Reserve-Ablehnung; Billing und Invoice
    nicht aufgerufen\\
  Invoice nicht verfügbar & Bestellung bleibt\newline
    \code{COMPLETED} & Versuche 1 und 2 eingeplant, Versuch 3 ausgeschöpft,
    Verarbeitungsfehler auditiert; kein Refund\\
  Warehouse-Commit fehlerhaft & HTTP 502,\newline
    \code{ROLLBACK_COMPLETED} & Refund über Fassade, Lagerfreigabe und je ein
    Snapshot für angeforderte und abgeschlossene Kompensationen\\
  \bottomrule
\end{tabularx}
\caption{End-to-End-Beobachtung im Präsentationsmodus}
\end{table}

\littleGap

Bei der Website-Prüfung des Warehouse-Fehlers wurden 22 chronologische
Snapshots unter derselben Correlation-ID aus PostgreSQL gelesen. Die Kette
enthielt unter anderem \code{STOCK_COMMIT_FAILED},
\code{PAYMENT_REFUND_REQUESTED}, \code{PAYMENT_REFUNDED},
\code{STOCK_RELEASE_REQUESTED}, \code{STOCK_RELEASED} und den terminalen
Bestellstatus. Gleichzeitig waren die providerbezogene Erstattung im
Billing-Log und Release-Aufrufe in Warehouse- und Backend-Logs auffindbar.

\midGap

\subsection{Bewertung und verbleibende Testlücken}

Die Tests decken die entscheidenden Kontrollflüsse ab, erreichen aber keine
Produktionsvalidierung, außerdem werden PostgreSQL und RabbitMQ in den
Shop-Integrationstests, durch kontrollierte Grenzen (MOCK's) ersetzt.\\
Softwareprojekte dieser Größe die länger betrieben und somit auch gepflegt werden müssen,
benötigen einen konsequenten Aufbau einer Testumgebung nach ISTQB\cite{istqb} und V-Modell.
Diese Umgebung wächst mit den Anforderungen des Projekts auf jeder Ebene. 
Aufgrund der verteilten Umgebung werden zusätzlich Last-, Stress- und Penetrationstests notwendig, 
auch \emph{Chaos-Engineering} und Partitions/Netzwerktests können zur Robustheit beitragen.

\midGap

\newpage

\section{Reflexion}
\label{sec:reflection}

\littleGap

\subsection{Gute Entscheidungen}

\begin{itemize}  
	\item \term{Zentrale Microservice-Architektur:}
    Alle wesentliche Zugriffe laufen über einen zentralen Shop-Service und dieser
		meldet nur zum StoreProxy durch. Die Architektur vermittelt eine gute Übersichtlichkeit
		und kann sowhol synchron als auch dynamisch asynchrone (Teil-)Prozesse durchführen und
		kann orchestriert werden.
  \item \term{Die Payment-Fassade stellt eine belastbare Grenze dar:}
    Kein Billing-Aufrufer greift direkt auf konkrete Adapter zu. Providerwahl,
    Timeout, Refund und Status sind einheitlich, und neue Adapter werden ohne
    Änderung bestehenden Fassadencodes entdeckt.
  \item \term{Fehler sind fachlich und technisch nachvollziehbar:}
    Correlation-ID, strukturierte Logs, OTLP-Daten und append-only
    Audit-Snapshots ergänzen sich. Auch Kompensationsschritte sind
    eigenständige Ereignisse.
  \item \term{Fault Injection:}
    Setzt an den tatsächlichen Servicegrenzen an. Sie
    simuliert keine bloßen Frontendmeldungen, sondern löst Refund, Release,
    Retry und DLQ-Verhalten aus.
  \item \term{Stubs für Payment-Provider, E-Mail und DB-Client sind bereits als Adapter bzw. Strategie entworfen:}
    Die entsprechenden Slots für diese Komponenten sind bereits gegeben, Kapselung und Erweiterbarkeit wurden berücksichtigt.
\end{itemize}

\midGap

\subsection{Weniger gute Entscheidungen}

\begin{enumerate}
  \item \term{Verwendung des \emph{TransactionalOutboxPattern}:} Zwischen erfolgreicher
    Zahlung und RabbitMQ-Publikation besteht eine \emph{Dual-Write-Grenze}.
    Publisher-Confirms und Refund reduzieren das Risiko, eine persistente
    Outbox im BillingService würde den Übergang jedoch noch robuster machen.  
  \item \term{Containerisierte Infrastruktur:} Die PowerShell-Skripte
    ermöglichen einen guten Windows-Demobetrieb, sind jedoch weniger portabel
    als Docker Compose oder eine containerbasierte Testumgebung. PostgreSQL,
    RabbitMQ und Collector sollten für CI reproduzierbar gekapselt werden. Dieses
		Feature konnte leider nicht mehr berücksichtigt werden, war aber geplant.
  \item \term{Infrastrukturtests früher einplanen:} Viele Fehlerlücken wurden
    erst durch manuelles Abschalten einzelner Services sichtbar. Frühe
    automatisierte Ausfalltests hätten Readiness, Timeout und Kompensation
    schneller zusammengeführt.
  \item \term{Retry-Richtlinie:} Drei Invoice-Versuche erfüllen zwar die 
		Anforderung, produktionstauglicher sind jedoch Techniken wie exponentielles Backoff mit Jitter,
		Metriken, Alarmierung und ein kontrollierter Redrive-Prozess für die DLQ.
\end{enumerate}

\midGap

\subsection{Eingegangene Kompromisse}

\begin{enumerate}
  \item \term{Hybrid aus gRPC und Events:} 
		Zwischen erfolgreicher synchroner Konsistenz für checkoutkritische Schritte, aber \emph{"`eventual consistency"'} bei Rechnung und 
		Auditing, nach Voraussetzung
  \item \term{Event Sourcing Light}
		Sehr gute Historie leider ohne vollständige Aggregate-Rekonstruktion; operative Daten bleiben separate Wahrheit		
  \item \term{Provider-Stubs} 
		Deterministische Tests ohne Geldbewegung; keine Aussage über Webhooks, Provider-SLAs.(War geplant)
  \item \term{Dateiablage für Lager und Rechnung:}
		Einfache, sichtbare Demonstration; keine horizontale Skalierung und eingeschränkte Parallelität, dafür gute Handhabbarkeit
  \item \term{Pickup-Outbox statt echter E-Mail} 
		Sichere und reproduzierbare Rechnungserzeugung; Versandzustellung wird nicht validiert, dafür vorbereitet mit \emph{Mailkit-Package}
  \item \term{Orchestrierte SAGA:} 
		Ablauf und Kompensation sind zentral verständlich; der ShopService trägt dafür wie geplant mehr Prozesswissen und stellt das Zentrum dar
	\item \term{Datenbank, E-Mail} 
		Die Komponenten funktionieren zwar nur eingeschränkt (DB-Cluster-Instanz, E-Mails werden nur in Ordner geschrieben/kein echter SMTP-Client); Aber die Komponenten sind leistungsfähig (QuestPDF/Mailkit) und der Datenbankadapter kapselt mit EF.Core
\end{enumerate}

\midGap


\subsection{Fazit}

Die Realisierung des Online-Stores mit aktueller Technologie und Systemen zeigt auf, wie stark sich grade die Webentwicklung in den letzten 20 Jahren verändert hat. Wo früher noch eine Website mit einem DOM und Model ausreichend war, sind inzwischen wesentliche Neuerungen Stand aktueller Technik, dies schlägt sich auch in der Projektgröße und Umfang nieder. Aufgrund der zahlreichen Anwendungsfälle die durch verteilte Systeme abgebildet werden können, steigt die Anzahl der Techniken, des notwendigen Konfigurationsmanagment und Modellwissens.\\
Wesentlich sind allerdings die Architekturmerkmale und das technische Verständnis hinsichtlich einer Microservice-Architektur. Auch wenn diese Umsetzung exemplarisch ist, war der Lerneffekt deutlich, die Realisierung eines Online-Stores "`State-of-the-Art"' stellt in der Tat eine gute Übung dar, um aktuelle Technologien, Topologien und Entwurfsmuster zu verwenden bzw. zu verstehen.\\

\midGap

\newpage

\section*{Literatur und Standards}
\addcontentsline{toc}{section}{Literatur und Standards}
\small
\begin{thebibliography}{9}

\bibitem{cgpt}
ChatGPT-Codex Agent, zur Erstellung des Web-Frontend's
\url{https://chatgpt.com/de-DE/codex/}

\bibitem{rfc9457}
IETF:
\emph{RFC 9457 - Problem Details for HTTP APIs}, 2023.
\url{https://www.rfc-editor.org/rfc/rfc9457}

\bibitem{openapi}
OpenAPI Initiative:
\emph{OpenAPI Specification 3.1}.
\url{https://spec.openapis.org/oas/v3.1.0}

\bibitem{gamma}
E. Gamma, R. Helm, R. Johnson und J. Vlissides:
\emph{Design Patterns: Elements of Reusable Object-Oriented Software}.
Addison-Wesley, 1994.

\bibitem{richardson}
C. Richardson
\emph{Microservices Patterns}. Manning, 2018,
Kapitel 4: Managing transactions with sagas.

\bibitem{istqb}
 International Software Testing Qualifications Board
\url{https://istqb.org}

\bibitem{ef}
EntityFrameworkCore
\url{https://learn.microsoft.com/de-de/ef/core/get-started/overview/first-app?tabs=netcore-cli}

\bibitem{yarp}
Microsofts 'Yet Another Reverse Proxy' für .NET
\url{https://dotnet.github.io/yarp/}

\bibitem{fowler}
M. Fowler:
\emph{Event Sourcing}, 2005.
\url{https://martinfowler.com/eaaDev/EventSourcing.html}

\bibitem{rabbitmq}
RabbitMQ:
\emph{Consumer Acknowledgements and Publisher Confirms}.
\url{https://www.rabbitmq.com/docs/confirms}

\bibitem{postgres}
PostgreSQL Global Development Group:
\emph{PostgreSQL 18 Documentation}.
\url{https://www.postgresql.org/docs/18/}

\end{thebibliography}

\end{document}
